\section{Related Work}
\label{sec:related}

\para{Personas and role prompts.} Assigning a role in the system prompt (``You are an expert physician'') is one of the most common prompting practices. \citet{zheng2023helpful} evaluate many social roles across models and factual question sets and find that adding a persona does not reliably improve accuracy; effects are small and hard to predict. \citet{salewski2023impersonation} show that impersonated personas can shift both task performance and biases. These studies focus on {\em short} role statements chosen to match the task. We instead study detailed, deployment-style personas (biography, tone, hobbies) that are unrelated to the task, and we measure output length as well as accuracy. Unlike prior work, we find that such personas roughly double output length.

\para{Format restrictions.} Production systems often require structured output. \citet{tam2024speak} report that format restrictions such as JSON can degrade reasoning, especially on math, and that looser formats hurt less. Building on this, we compare three format prompts (a final-answer line, markdown sections and JSON) against content-free controls, and we score each response under both a lenient and a strict parser. This lets us show that the same format prompt both lowers content accuracy and raises parser accuracy, and that a format rule written for one task type can misfire on another.

\para{Over-refusal and safety behaviour.} Safety training can cause models to refuse benign requests. XSTest~\cite{rottger2023xstest} and OR-Bench~\cite{cui2024orbench} build prompts that look harmful but are safe, and measure how often models refuse them. These benchmarks vary the {\em user request}. We vary the {\em system prompt} and keep ordinary academic questions fixed, which asks a different question: does safety text in the system prompt spill over onto requests that are clearly benign? For the models we study, it mostly does not.

\para{Prompt sensitivity and evaluation.} Accuracy can change by large amounts under superficial prompt-format changes~\cite{sclar2023quantifying}, and single-prompt evaluations can give misleading rankings~\cite{mizrahi2024state}. Our work shares the concern that evaluation conditions differ from deployment conditions, and adds a specific, practical axis of variation: the categories of system prompt that deployments actually use. Following this line of work, we use several wording variants per category and report variant-level heterogeneity.

\para{Reasoning length.} Chain-of-thought prompting improves multi-step reasoning~\cite{wei2022chain,kojima2022large}. Several of the system prompts we study shorten or remove visible reasoning as a side effect (brevity limits, ``do not show your work'', compact JSON fields). Our results show that these side effects are the main source of accuracy loss, which is consistent with the chain-of-thought literature but arises from prompts that were not written to affect reasoning at all.
